\section{문제, 현황, 그리고 이 보고서의 위치}

\subsection{문제와 표기}
$\R^n$ 의 격자 $L$ (covolume 1)에 대해
\[
V(L):=\Vol\{\text{원점 중심의, $L\setminus\{0\}$ 과 만나지 않는 가장 큰 열린 공}\},\qquad V_n:=\sup_L V(L)
\]
로 두면 격자 구 채우기 밀도는 $\delta(L)=V(L)/2^n$, $\delta_n=V_n/2^n$ 이다. 동치로, 고정된 격자 $L$ 과 양의 정부호 행렬 $A$ 에 대해 타원체 $\mathcal E_A=\{x: Ax\cdot x<1\}$ 가 $L$-free ($\mathcal E_A\cap L=\{0\}$) 이면 $V\ge \Vol(\mathcal E_A)/\covol(L)$ 이다. 이 보고서 전체에서 $N:=n(n+1)/2=\dim\Sym_n$ 이다.

\subsection{알려진 결과 (2026년 11월 기준, 1회차 문헌 조사로 확정)}
\begin{center}
\begin{tabular}{@{}llll@{}}
\toprule
결과 & $V_n$ 에 대한 한계 & 차원 & 방법 \\
\midrule
Minkowski--Hlawka & $\ge 2\zeta(n)$ & 모든 $n$ & 평균(Siegel) \\
Rogers 1947, Ball 1992 & $\ge 2(n-1)\zeta(n)$ & 모든 $n$ & 평균 + successive minima \\
Vance 2011 & $\ge (6/e)\,n$ & $4\mid n$ & Hurwitz 사원수 가군 \\
Venkatesh 2013 & $\ge \tfrac12 n\log\log n$ & 희소 & 원분체 $\Z[\zeta_m]$ 가군, 궤도 크기 $m$ \\
Gargava--Viazovska 2024 & $\ge n\log\log n$ & 희소 & 랜덤 ideal 격자 \\
Campos--Jenssen--Michelen--Sahasrabudhe & $\ge(\tfrac12-o(1))n\log n$ & 모든 $n$ (비격자) & 그래프 이론 \\
\textbf{Klartag 2025} & $\ge c\,n^2$ & 모든 $n$ & 확률적으로 진화하는 타원체 \\
Rezvyakova 2026 & $\ge(\tfrac1{2e}-o(1))\,n^2$ & $n\to\infty$ & Klartag 증명의 상수 추적 \\
Abuya--Gargava--Zhao 2026 & $\ge c\,n^2\log\log n$ & 희소 & Klartag $\times$ 원분 대칭 \\
\textbf{Gargava--Gr\"obner--Simkin 2026} & $\ge c\,n^2\log\log n$ & 모든 큰 $n$ & 원분 접착(gluing) 구성 \\
\midrule
Kabatiansky--Levenshtein 1978 & $\le 2^{0.401n}$ & & 구면 부호 \\
OpenAI ``Ten Advances'' 2026 (Ch.\,1--2) & $\le 2^{(0.3956+o(1))n}$ & & Cohn--Elkies LP 의 정확한 지수 \\
Venkatesh 의 추측 & $\le n^2(\log n)^{O(1)}$ ? & & ``격자의 자유도는 $n^2$'' \\
\bottomrule
\end{tabular}
\end{center}
사용자가 올려 준 두 논문은 위 표의 Klartag 2025(하한)와 OpenAI 2026(상한)이며, ``$\log\log n$ 이 붙은 최신 결과''는 Abuya--Gargava--Zhao (arXiv:2606.05105) 와 그것을 모든 차원으로 확장한 Gargava--Gr\"obner--Simkin (arXiv:2610.12619) 이다.

\subsection{Klartag 방법의 핵심 (이 보고서에서 계속 쓰는 언어)}
격자 $L$ 을 고정하면 $L$-free 행렬 전체
\[
\mathcal R_L=\{A\in\Sym_n:\ Ax\cdot x\ge 1\ \ \forall x\in L\setminus\{0\}\}
\]
는 국소 유한 다면체(\emph{Ryshkov 다면체})이고, 꼭짓점이 Voronoi 의 \emph{perfect form}, $\det$ 의 국소 최소점이 \emph{extreme form} 이다. Klartag 과정은 $\mathcal R_L$ 안에서의 브라운 운동 $dA_t=\pi_t(dW_t)$ 으로, 면에 닿으면(=격자점이 타원체 경계에 닿으면) 그 면에 달라붙고(sticky), 자유도 $N_t=\dim F_t$ 가 0이 되는 꼭짓점에서 멈춘다. $\log\det$ 의 오목성으로
$\ \E\log\det A_T\le n\log a_0-\tfrac12\int_0^T\E[\|A_t\|_{op}^{-2}N_t]\,dt$
이고, 접촉점 수의 기댓값은 반사원리와 Siegel 평균값 공식(1차 모멘트)만으로 $\lesssim e^{n^2t/8}$ 로 제어된다. 요점은 \emph{``자유도 하나가 접촉쌍 하나를 흡수한다''} 는 회계이며, 대칭 격자 버전(AGZ, GGS)은 \emph{``궤도 하나가 자유도 하나만 소모한다''} 로 이 회계를 개선한 것이다:
\[
V\ \gtrsim\ \frac{(\text{궤도 크기 } w)\times(\text{불변 이차형식 공간의 차원 } D)}{2e}.
\]

\subsection{이 연구 세션이 한 일과 하지 못한 일 (정직한 요약)}
\begin{itemize}
\item \textbf{하지 못한 일.} $c\,n^2\log\log n\cdot 2^{-n}$ 의 \emph{차수}를 개선하는 정리는 얻지 못했다. 이는 오래된 미해결 문제이며, 5시간짜리 탐색의 현실적인 목표는 (i) 방법의 한계를 정확히 규명하고 (ii) 검증 가능한 부분 결과와 (iii) 구체적인 다음 공격 지점을 남기는 것이었다.
\item \textbf{한 일.} \S\ref{sec:summary} 의 표에 [증명됨]/[조건부]/[추측]/[수치관찰] 로 구분해 정리했다. 요지는 다음 다섯 가지다.
 \begin{enumerate}
 \item \textbf{(통일과 한계)} 대칭을 이용하는 모든 알려진 논증을 하나의 이득 공식 $g=|G|/(2\dim_\Q A)=\tfrac12\sum_{g\in G}\rho(g)^2$ 로 통일하고, \emph{고정점 없는 작용}으로는 $\log\log$ 를 넘을 수 없음을 증명했다. 비자유 작용에 대해서는 ``텐서 안정성'' 이라는 판정 기준과 여러 no-go 정리를 얻었다 (\S\ref{sec:sym}).
 \item \textbf{(명시적 상수)} Hurwitz/icosian 가군 격자에 대해 Klartag 형 한계 $V_n\ge(\tfrac3{2e}-o(1))n^2$ ($4\mid n$), $(\tfrac{15}{4e}-o(1))n^2$ ($8\mid n$) 를 완전한 증명으로 쓰고 독립 심사를 거쳤다 --- Rezvyakova 상수 $\tfrac1{2e}$ 의 3배, 7.5배. 차수는 $n^2$ 그대로다 (\S\ref{sec:quat}).
 \item \textbf{(유일하게 남은 차수 개선 후보)} 자유 작용의 한계를 피하는 유일한 후보로 $SL_2(p)$ Weil 격자의 텐서족을 찾았고, $p=13$ (Leech 격자)에서 ``정확히 경계 안정'' 임을 계산기 보조 정리로 증명했다. 일반 $p$ 에서의 성립 여부는 Weil 격자(=곡선 $y^2=x^p-x$ 의 Mordell--Weil 격자)의 산술 문제로 환원되며, 성립한다면 $n^{2+c}$ 형의 하한을 주게 되므로(기존 추측과 충돌) 어느 쪽이든 흥미롭다 (\S\ref{sec:leech}).
 \item \textbf{(모형)} Klartag 과정을 랜덤 선형계획(음의 마진 퍼셉트론)으로 재해석하는 Poisson 모형을 세워 상수의 손실($\tfrac1{2e}\to\tfrac12\to1$)을 설명하고, 모형에서의 정확한 법칙 $V_{LP}\sim\chi^2_{2N}$ 을 얻었다. 그러나 실제 랜덤 격자에서는 $n\le28$ 까지 이 예측이 성립하지 않음도 실험으로 확인했다 (\S\ref{sec:perceptron}).
 \item \textbf{(소차원의 직접 구성)} $n\le12$ 의 비대칭 탐색과 $n\le24$ 의 대칭(등변) 탐색을 구현·검증했다. 대칭 과정은 $E_8,K_{12},BW_{16},\Lambda_{20}$, Leech 를 실제로 찾아내며, ``$V\approx$ (궤도 크기/2)$\times$(불변형식 차원)'' 이라는 스케일링 법칙이 관찰된다 (\S\ref{sec:numerics}, \S\ref{sec:equiv}).
 \end{enumerate}
\end{itemize}

\subsection{연구 절차}
요청받은 절차(로그 읽기 $\to$ 과제 선택 $\to$ 증명 시도/반례 탐색/계산 검증 $\to$ 막힌 지점 기록 $\to$ 로그 저장)를 회차 단위로 반복했다. 각 회차는 ``증명/계산 담당'' 하위 에이전트와, 그 결과를 \emph{반박하도록} 지시받은 독립 ``심사'' 에이전트의 쌍으로 구성했고, 심사를 통과하지 못한 주장은 강등하거나 폐기했다. 모든 수치는 실제 실행한 코드에서 나온 것이며 코드와 원자료는 저장소의 \texttt{sphere\_packing\_research/rounds/} 아래에 있다. 회차별 기록은 \texttt{research\_log.md} 에 있다.
